\documentclass[11pt,a4paper]{ctexart}
\usepackage{amsmath,amssymb,bm,mathtools,geometry,enumitem,booktabs,xcolor,hyperref,fancyhdr,microtype}
\geometry{left=25mm,right=25mm,top=23mm,bottom=23mm,headheight=15pt}
\definecolor{ink}{HTML}{18334B}\definecolor{accent}{HTML}{176B77}
\hypersetup{colorlinks=true,linkcolor=ink,urlcolor=accent}
\pagestyle{fancy}\fancyhf{}\fancyhead[L]{\small MVDR 与原问题的最优性}\fancyhead[R]{\small 理论分析}\fancyfoot[C]{\thepage}\renewcommand{\headrulewidth}{0.3pt}
\setlength{\parindent}{2em}\setlength{\parskip}{3pt}\setlist{itemsep=2pt,topsep=3pt}
\allowdisplaybreaks
\title{\textbf{MVDR 接收波束的最优性证明}\\[4pt]\large 基于多 UAV 辅助 ISCC 论文原式的条件与证明}
\author{多 UAV 辅助 ISCC 系统}
\date{2026 年 10 月 28 日}
\begin{document}\maketitle
\noindent\textbf{结论。}在第\,三\,节列出的模型条件下，\textbf{MVDR 接收波束可以作为系统能耗最小化问题的最优接收波束}。如果原问题存在全局最优解，那么至少有一个全局最优方案使用 MVDR 接收波束；在优化其他变量时，根据对应的位置、匹配和功率计算 MVDR 接收波束，仍能达到原问题的最低系统能耗。

证明思路是：先求出使感知信干噪比最大的接收波束，再证明通过调整感知发射波束，可以保持原有感知性能，同时不增加发射能耗，也不降低通信性能。由此得到使用 MVDR 接收波束的最优方案。

\noindent\textbf{符号约定。}下文沿用系统模型的原符号。上标 $H$ 表示共轭转置，$\bm h_{c_j,u_i}(t)$ 为 $1\times N$ 行向量。推导中的 $t$ 和 $i$ 为任意时隙和任意 UAV，结论对所有时隙和 UAV 分别适用。
\begingroup\small\tableofcontents\endgroup
\newpage

\section{论文中接收波束影响哪些量}
\subsection{问题 P0 的完整形式}
\begingroup\small
\begin{align}
    \mathrm{P0}:&
    \notag\\
    \min&
    \frac{1}{t_{\max}}
    \sum_{t=1}^{t_{\max}}
    \Bigg[
    \omega_{1}\kappa
    \Bigg(
        \sum_{u_i\in\mathcal{U}}
        C^{\mathrm{sen}}\bar{D}^{\mathrm{sen}}
        R_i^{\mathrm{sen}}(t)
        \left(f_{u_i}(t)\right)^2
        +\sum_{c_j\in\mathcal{C}}
        C_j(t)L_j(t)
        \left(f_{c_j}(t)\right)^2
    \Bigg)
    \notag\\
    &
    +\omega_{2}
    \sum_{u_i\in\mathcal{U}}
    \Bigg(
        \bar{D}^{\mathrm{sen}}
        \left\|\mathbf{w}_i(t)\right\|^2
        +D_{u_i}^{\mathrm{off}}(t)
        \left\|\mathbf{b}_i(t)\right\|^2
        +\frac{
            \varrho_{1}
            \left\|
                \mathbf{q}_{u_i}(t+1)-\mathbf{q}_{u_i}(t)
            \right\|^3
        }{\tau^2}
        +\frac{
            \varrho_{2}\tau^2
        }{
            \left\|
                \mathbf{q}_{u_i}(t+1)-\mathbf{q}_{u_i}(t)
            \right\|
        }
    \Bigg)
    \notag\\
    &
    +\omega_{3}
    \sum_{c_j\in\mathcal{C}}
    D_{c_j}^{\mathrm{off}}(t)p_j(t)
    \Bigg]
    \notag\\
    %
    \text{s.t.}&
    \sum_{j=1}^{J}\eta_{i,j}(t)=1,
    \forall u_i\in\mathcal{U},
    \forall t\in\mathcal{T},
    \notag\\
    %
    &
    \sum_{i=1}^{I}\eta_{i,j}(t)\leq1,
    \forall c_j\in\mathcal{C},
    \forall t\in\mathcal{T},
    \notag\\
    %
    &
    V_{\min}\tau
    \leq
    \left\|
        \mathbf{q}_{u_i}(t+1)-\mathbf{q}_{u_i}(t)
    \right\|
    \leq V_{\max}\tau,
    \forall u_i\in\mathcal{U},
    \forall t\in\mathcal{T},
    \notag\\
    %
    &
    \left\|
        \mathbf{q}_{u_{i_1}}(t)-\mathbf{q}_{u_{i_2}}(t)
    \right\|
    \geq d_{\min}^{\mathrm{UAV}},
    \forall u_{i_1},u_{i_2}\in\mathcal{U},
    u_{i_1}\neq u_{i_2},
    \forall t\in\mathcal{T},
    \notag\\
    %
    &
    \left\|\mathbf{w}_i(t)\right\|^2
    \leq P_{\mathrm{UAV}}^{\max},
    \forall u_i\in\mathcal{U},
    \forall t\in\mathcal{T},
    \notag\\
    %
    &
    \left\|\mathbf{g}_i(t)\right\|^2=1,
    \forall u_i\in\mathcal{U},
    \forall t\in\mathcal{T},
    \notag\\
    %
    &
    \frac{
        \left|
            \mathbf{g}_i^H(t)
            \mathbf{A}(\theta_i(t))
            \mathbf{w}_i(t)
        \right|^2
    }{
        \displaystyle
        \sum_{j=1}^{J}
        \eta_{i,j}(t)p_j(t)
        \left|
            \mathbf{g}_i^H(t)
            \mathbf{h}_{c_j,u_i}^H(t)
        \right|^2
        +\sigma^2
    }
    \geq\epsilon,
    \forall u_i\in\mathcal{U},
    \forall t\in\mathcal{T},
    \notag\\
    %
    &
    \left\|\mathbf{b}_i(t)\right\|^2
    \leq P_{\mathrm{UAV}}^{\max},
    \forall u_i\in\mathcal{U},
    \forall t\in\mathcal{T},
    \notag\\
    %
    &
    D_{u_i}^{\mathrm{off}}(t)
    R_i^{\mathrm{off}}(t)
    \geq
    \bar{D}^{\mathrm{sen}}
    R_i^{\mathrm{sen}}(t),
    \forall u_i\in\mathcal{U},
    \forall t\in\mathcal{T},
    \notag\\
    %
    &
    \bar{D}^{\mathrm{sen}}
    +D_{u_i}^{\mathrm{off}}(t)
    \leq
    \sum_{j=1}^{J}
    \eta_{i,j}(t)D_{c_j}^{\mathrm{off}}(t),
    \forall u_i\in\mathcal{U},
    \forall t\in\mathcal{T},
    \notag\\
    %
    &
    \sum_{i=1}^{I}
    \eta_{i,j}(t)
    \bar{D}^{\mathrm{sen}}
    R_{i,j}^{(1)}(t)
    +\sum_{i=1}^{I}
    \eta_{i,j}(t)
    D_{u_i}^{\mathrm{off}}(t)
    R_{i,j}^{(2)}(t)
    \notag\\
    &
    +\Bigg(
        D_{c_j}^{\mathrm{off}}(t)
        -\sum_{i=1}^{I}
        \eta_{i,j}(t)\bar{D}^{\mathrm{sen}}
        -\sum_{i=1}^{I}
        \eta_{i,j}(t)D_{u_i}^{\mathrm{off}}(t)
    \Bigg)
    R_j^{(3)}(t)
    \geq L_j(t),
    \forall c_j\in\mathcal{C},
    \forall t\in\mathcal{T},
    \notag\\
    %
    &
    p_j(t)\leq P_{\mathrm{CU}}^{\max},
    \forall c_j\in\mathcal{C},
    \forall t\in\mathcal{T},
    \notag\\
    %
    &
    \bar{D}^{\mathrm{sen}}
    +D_{u_i}^{\mathrm{off}}(t)
    +\frac{
        C^{\mathrm{sen}}
        \bar{D}^{\mathrm{sen}}
        R_i^{\mathrm{sen}}(t)
    }{
        f_{u_i}(t)
    }
    \leq D_{\max}^{\mathrm{sen}},
    \forall u_i\in\mathcal{U},
    \forall t\in\mathcal{T},
    \notag\\
    %
    &
    D_{c_j}^{\mathrm{off}}(t)
    +\frac{
        C_j(t)L_j(t)
    }{
        f_{c_j}(t)
    }
    \leq D_j^{\max}(t),
    \forall c_j\in\mathcal{C},
    \forall t\in\mathcal{T},
    \notag\\
    %
    &
    \sum_{u_i\in\mathcal{U}}f_{u_i}(t)
    +\sum_{c_j\in\mathcal{C}}f_{c_j}(t)
    \leq F_{\max},
    \forall t\in\mathcal{T},
    \notag\\
    %
    \text{var.}&
    \mathbf{w}_i(t)\geq0,
    \mathbf{b}_i(t)\geq0,
    \mathbf{g}_i(t)\geq0,
    \mathbf{q}_{u_i}(t),
    D_{u_i}^{\mathrm{off}}(t)>0,
    f_{u_i}(t)>0,
    \forall u_i\in\mathcal{U},
    \forall t\in\mathcal{T},
    \notag\\
    %
    &
    p_j(t)>0,
    D_{c_j}^{\mathrm{off}}(t)>0,
    f_{c_j}(t)>0,
    \forall c_j\in\mathcal{C},
    \forall t\in\mathcal{T},
    \eta_{i,j}(t)\in\{0,1\},
    \forall u_i\in\mathcal{U},
    \forall c_j\in\mathcal{C},
    \forall t\in\mathcal{T}.
    \notag
\end{align}
\endgroup

\subsection{接收波束对约束和目标的影响}
接收波束 $\bm g_i(t)$ 在 $P0$ 中直接受到单位范数约束
\begin{equation}\label{eq:norm}
\|\bm g_i(t)\|^2=1.
\end{equation}
在感知 SINR 约束中，$\bm g_i(t)$ 同时影响分子中的目标回波增益与分母中的 CU 干扰功率：
\begin{equation}\label{eq:sinr}
\frac{\left|\bm g_i^H(t)\mathbf A(\theta_i(t))\bm w_i(t)\right|^2}
{\displaystyle\sum_{j=1}^{J}\eta_{i,j}(t)p_j(t)
\left|\bm g_i^H(t)\bm h_{c_j,u_i}^{H}(t)\right|^2+\sigma^2}\ge\epsilon.
\end{equation}
其中，论文采用的感知信道为
\begin{equation}\label{eq:A}
\mathbf A(\theta_i(t))=\sqrt{\xi_0\rho d_i^{-4}(t)}\,
\bm a_r(\theta_i(t))\bm a_t^H(\theta_i(t)).
\end{equation}
感知信息速率由同一感知 SINR 决定：
\begin{equation}\label{eq:rate}
R_i^{\rm sen}(t)=\frac{\delta}{2\nu}\log_2\!\left(
1+\frac{2\sigma_{\rm pre}^2\gamma^2B^3\nu
\left|\bm g_i^H(t)\mathbf A(\theta_i(t))\bm w_i(t)\right|^2}
{\displaystyle\sum_{j=1}^{J}\eta_{i,j}(t)p_j(t)
\left|\bm g_i^H(t)\bm h_{c_j,u_i}^{H}(t)\right|^2+\sigma^2}\right).
\end{equation}
因此，$\bm g_i(t)$ 还通过 $R_i^{\rm sen}(t)$ 间接影响感知任务的卸载约束与计算时延约束：
\begin{align}
D_{u_i}^{\rm off}(t)R_i^{\rm off}(t)&\ge\bar D^{\rm sen}R_i^{\rm sen}(t),\label{eq:impact-off}\\
\bar D^{\rm sen}+D_{u_i}^{\rm off}(t)+
\frac{C^{\rm sen}\bar D^{\rm sen}R_i^{\rm sen}(t)}{f_{u_i}(t)}
&\le D_{\max}^{\rm sen}.\label{eq:impact-delay}
\end{align}
原因是感知数据量为 $\bar D^{\rm sen}R_i^{\rm sen}(t)$：它决定需要卸载和计算的数据量，也出现在 $P0$ 目标函数的 BS 感知计算能耗中。其他约束在其余变量固定时不直接含 $\bm g_i(t)$；在联合优化时，卸载、计算和功率等变量仍会相互耦合。

\section{MVDR 为什么不依赖尚未求得的发射波束}
\subsection{代入论文原式，不省略物理变量}
由式\,(\ref{eq:A}) 直接得到
\begin{align}
&\left|\bm g_i^H(t)\mathbf A(\theta_i(t))\bm w_i(t)\right|^2\notag\\
&=\xi_0\rho d_i^{-4}(t)
\left|\bm g_i^H(t)\bm a_r(\theta_i(t))\right|^2
\left|\bm a_t^H(\theta_i(t))\bm w_i(t)\right|^2.\label{eq:factor}
\end{align}
因此原 SINR 完整地等于
\begin{equation}\label{eq:factored}
\frac{\xi_0\rho d_i^{-4}(t)
\left|\bm a_t^H(\theta_i(t))\bm w_i(t)\right|^2
\left|\bm g_i^H(t)\bm a_r(\theta_i(t))\right|^2}
{\displaystyle\sum_{j=1}^J\eta_{i,j}(t)p_j(t)
\left|\bm g_i^H(t)\bm h_{c_j,u_i}^{H}(t)\right|^2+\sigma^2}.
\end{equation}
考虑固定 UAV 位置、频谱共享变量 $\eta_{i,j}(t)$、CU 功率 $p_j(t)$ 和感知发射波束 $\bm w_i(t)$ 的情况。式\,(\ref{eq:factored}) 的分子中，下面这一项表示发射波束指向目标的增益及感知信道的路径增益：
\begin{equation}\label{eq:tx-factor}
\xi_0\rho d_i^{-4}(t)
\left|\bm a_t^H(\theta_i(t))\bm w_i(t)\right|^2.
\end{equation}
这一项只包含目标参数、位置及发射波束，没有接收波束 $\bm g_i(t)$；所以改变 $\bm g_i(t)$ 时，式\,(\ref{eq:tx-factor}) 的值不变。若该值大于零，将其作为正因子从最大化问题中提出，不改变最大值对应的接收波束。因此，求解接收波束时只需最大化：
\begin{equation}\label{eq:receive-optimization}
\max_{\|\bm g_i(t)\|^2=1}
\frac{\left|\bm g_i^H(t)\bm a_r(\theta_i(t))\right|^2}
{\displaystyle\sum_{j=1}^J\eta_{i,j}(t)p_j(t)
\left|\bm g_i^H(t)\bm h_{c_j,u_i}^{H}(t)\right|^2+\sigma^2}.
\end{equation}
式\,(\ref{eq:receive-optimization}) 中不再出现 $\bm w_i(t)$，这说明最优接收方向由目标接收导向、CU 信道、匹配、功率及噪声决定，与发射波束的具体取值无关。若式\,(\ref{eq:tx-factor}) 为零，所有接收波束的感知 SINR 均为零，另见本节的退化情形。

\subsection{最大 SINR 方向及严格上界}
为求解式\,(\ref{eq:receive-optimization})，先将分母中的干扰和噪声统一写成关于 $\bm g_i(t)$ 的二次型。由于接收波束满足 $\|\bm g_i(t)\|^2=1$，所以
\begin{equation}
\sigma^2=\sigma^2\bm g_i^H(t)\bm g_i(t)
=\bm g_i^H(t)(\sigma^2\mathbf I)\bm g_i(t).
\end{equation}
将噪声和每个 CU 的干扰一起写入括号，原分母等于
\begin{align}\label{eq:denominator}
&\sum_{j=1}^J\eta_{i,j}(t)p_j(t)
\left|\bm g_i^H(t)\bm h_{c_j,u_i}^{H}(t)\right|^2+\sigma^2\notag\\
&=\bm g_i^H(t)\!\left(\sum_{j=1}^J\eta_{i,j}(t)p_j(t)
\bm h_{c_j,u_i}^{H}(t)\bm h_{c_j,u_i}(t)+\sigma^2\mathbf I\right)\!\bm g_i(t).
\end{align}
假设 $\sigma^2>0$ 且 $\eta_{i,j}(t)p_j(t)\ge0$。对任意非零向量 $\bm v\in\mathbb C^{N\times1}$，有
\begin{align}
&\bm v^H\!\left(\sum_{j=1}^J\eta_{i,j}(t)p_j(t)
\bm h_{c_j,u_i}^{H}(t)\bm h_{c_j,u_i}(t)+\sigma^2\mathbf I\right)\!\bm v\notag\\
&=\sum_{j=1}^J\eta_{i,j}(t)p_j(t)
\left|\bm h_{c_j,u_i}(t)\bm v\right|^2+\sigma^2\|\bm v\|^2>0.
\end{align}
\begin{samepage}
因此，式\,(\ref{eq:denominator}) 括号中的矩阵为 Hermitian 正定矩阵，其正定平方根及逆平方根均存在。为了把分子和这个分母联系起来，在分子的内积中插入该矩阵的平方根及逆平方根，得到
\begin{align}\label{eq:sqrt-insertion}
&\bm g_i^H(t)\bm a_r(\theta_i(t))\notag\\
&=\left[\left(\sum_{j=1}^J\eta_{i,j}(t)p_j(t)
\bm h_{c_j,u_i}^{H}(t)\bm h_{c_j,u_i}(t)+\sigma^2\mathbf I\right)^{\frac12}\!\bm g_i(t)\right]^H\notag\\
&\quad\times\left[\left(\sum_{j=1}^J\eta_{i,j}(t)p_j(t)
\bm h_{c_j,u_i}^{H}(t)\bm h_{c_j,u_i}(t)+\sigma^2\mathbf I\right)^{-\frac12}\!\bm a_r(\theta_i(t))\right].
\end{align}
\end{samepage}
\begin{samepage}
这里两个矩阵相乘为单位矩阵。对式\,(\ref{eq:sqrt-insertion}) 右侧的两个向量应用复数 Cauchy--Schwarz 不等式 $|\bm x^H\bm y|^2\le\|\bm x\|^2\|\bm y\|^2$，得到
\begin{align}\label{eq:cs-norms}
&\left|\bm g_i^H(t)\bm a_r(\theta_i(t))\right|^2\notag\\
&\le\left\|\left(\sum_{j=1}^J\eta_{i,j}(t)p_j(t)
\bm h_{c_j,u_i}^{H}(t)\bm h_{c_j,u_i}(t)+\sigma^2\mathbf I\right)^{\frac12}\!\bm g_i(t)\right\|^2\notag\\
&\quad\times\left\|\left(\sum_{j=1}^J\eta_{i,j}(t)p_j(t)
\bm h_{c_j,u_i}^{H}(t)\bm h_{c_j,u_i}(t)+\sigma^2\mathbf I\right)^{-\frac12}\!\bm a_r(\theta_i(t))\right\|^2.
\end{align}
\end{samepage}
将两个范数平方分别展开，就得到：
\begin{align}
&\left|\bm g_i^H(t)\bm a_r(\theta_i(t))\right|^2\notag\\
&\le\bm g_i^H(t)\!\left(\sum_{j=1}^J\eta_{i,j}(t)p_j(t)
\bm h_{c_j,u_i}^{H}(t)\bm h_{c_j,u_i}(t)+\sigma^2\mathbf I\right)\!\bm g_i(t)\notag\\
&\quad\times\bm a_r^H(\theta_i(t))\!\left(\sum_{j=1}^J\eta_{i,j}(t)p_j(t)
\bm h_{c_j,u_i}^{H}(t)\bm h_{c_j,u_i}(t)+\sigma^2\mathbf I\right)^{-1}\!\bm a_r(\theta_i(t)).\label{eq:cs}
\end{align}
式\,(\ref{eq:cs}) 右侧第一项恰好是式\,(\ref{eq:denominator}) 的分母；将式\,(\ref{eq:cs}) 两侧除以该正数，就得到式\,(\ref{eq:receive-optimization}) 的上界。Cauchy--Schwarz 不等式取等当且仅当式\,(\ref{eq:sqrt-insertion}) 中的两个向量线性相关；将这一条件两侧乘以对应的逆平方根，可得 $\bm g_i(t)$ 与下式同方向，即相差非零复数倍：
\begin{equation}\label{eq:direction}
\left(\sum_{j=1}^J\eta_{i,j}(t)p_j(t)
\bm h_{c_j,u_i}^{H}(t)\bm h_{c_j,u_i}(t)+\sigma^2\mathbf I\right)^{-1}
\bm a_r(\theta_i(t)).
\end{equation}
由此，满足论文单位范数要求的一个解为
\begin{equation}\label{eq:mvdr}
\bm g_i^{\rm MVDR}(t)=
\frac{\left(\displaystyle\sum_{j=1}^J\eta_{i,j}(t)p_j(t)
\bm h_{c_j,u_i}^{H}(t)\bm h_{c_j,u_i}(t)+\sigma^2\mathbf I\right)^{-1}
\bm a_r(\theta_i(t))}
{\left\|\left(\displaystyle\sum_{j=1}^J\eta_{i,j}(t)p_j(t)
\bm h_{c_j,u_i}^{H}(t)\bm h_{c_j,u_i}(t)+\sigma^2\mathbf I\right)^{-1}
\bm a_r(\theta_i(t))\right\|}.
\end{equation}
其任意整体相位旋转也达到相同 SINR。经典无失真 MVDR 采用 $\bm g_i^H(t)\bm a_r(\theta_i(t))=1$ 的幅度归一化；式\,(\ref{eq:mvdr}) 采用论文要求的单位范数，方向相同，约束不同。Capon [1] 是相关的原始文献；Vorobyov [2, 第 2 节] 讨论最大 SINR、MVDR 和模型误差。

\textbf{退化情形。}若式\,(\ref{eq:tx-factor}) 为零，包括 $\bm a_t^H(\theta_i(t))\bm w_i(t)=0$ 的情形，则回波分子为零，所有接收方向的 SINR 都为零；在 $\epsilon>0$ 时该方案不可行。上面的唯一方向结论针对目标增益非零、目标导向矢量非零的情况。

\section{MVDR 接收波束的全局最优性证明}

\noindent\textbf{定理。}考虑论文中的系统能耗最小化问题 $P0$。假设 $P0$ 存在全局最优解，感知信道矩阵的完整表达式为
\begin{equation}\label{eq:theorem-channel}
\mathbf A(\theta_i(t))=\sqrt{\xi_0\rho d_i^{-4}(t)}\,
\bm a_r(\theta_i(t))\bm a_t^H(\theta_i(t)).
\end{equation}
该矩阵由接收与发射导向向量的外积构成，其每一列都是接收导向向量的标量倍数，故其秩不超过 $1$；当 $\xi_0\rho>0$ 且两个导向向量均非零时，矩阵非零，故其秩为 $1$。对所有 UAV 和时隙，设 $\sigma^2>0$、$p_j(t)\ge0$、$\bm a_r(\theta_i(t))\ne\bm0$、$\omega_2\ge0$，感知发射、卸载发射和接收波束均属于 $\mathbb C^{N\times1}$。

则 $P0$ 至少存在一个全局最优解，其每个 UAV 在每个时隙的接收波束可以取为
\begin{equation}\label{eq:unified-theorem}
\bm g_i^{\rm MVDR}(t)=
\frac{\left(\displaystyle\sum_{j=1}^J\eta_{i,j}(t)p_j(t)
\bm h_{c_j,u_i}^{H}(t)\bm h_{c_j,u_i}(t)+\sigma^2\mathbf I\right)^{-1}
\bm a_r(\theta_i(t))}
{\left\|\left(\displaystyle\sum_{j=1}^J\eta_{i,j}(t)p_j(t)
\bm h_{c_j,u_i}^{H}(t)\bm h_{c_j,u_i}(t)+\sigma^2\mathbf I\right)^{-1}
\bm a_r(\theta_i(t))\right\|}.
\end{equation}

\medskip
\noindent\textbf{证明。}分两步证明。首先从感知信道及 SINR 的原始表达式出发，推导式~\eqref{eq:unified-theorem}；然后构造可行解，证明该接收波束可用于 $P0$ 的全局最优解。

\subsection{第一步：推导 MVDR 接收波束}

任取 UAV $u_i$ 和时隙 $t$，固定 $P0$ 中除 $\bm g_i(t)$ 以外的所有优化变量，即各时隙的
\begin{align*}
&\bm w_{i'}(t'),\ \bm b_{i'}(t'),\ \bm q_{u_{i'}}(t'),\ D_{u_{i'}}^{\rm off}(t'),\ f_{u_{i'}}(t'),\\
&p_{j'}(t'),\ D_{c_{j'}}^{\rm off}(t'),\ f_{c_{j'}}(t'),\ \eta_{i',j'}(t'),\\
&\bm g_{i'}(t')\quad\text{其中 }(i',t')\ne(i,t).
\end{align*}
以上指标 $i',j',t'$ 遍历各 UAV、CU 和时隙。考虑以下原始感知信道模型和感知 SINR：
\begin{equation}\label{detail:eq:channel}
\mathbf A(\theta_i(t))
=\sqrt{\xi_0\rho d_i^{-4}(t)}\,
\bm a_r(\theta_i(t))\bm a_t^H(\theta_i(t)),
\end{equation}
\begin{equation}\label{detail:eq:original-sinr}
\operatorname{SINR}_i(t)=
\frac{\left|\bm g_i^H(t)\mathbf A(\theta_i(t))\bm w_i(t)\right|^2}
{\displaystyle\sum_{j=1}^J\eta_{i,j}(t)p_j(t)
\left|\bm g_i^H(t)\bm h_{c_j,u_i}^{H}(t)\right|^2+\sigma^2}.
\end{equation}
其中 $\bm a_r$、$\bm a_t$ 是接收和发射导向向量，$\bm h_{c_j,u_i}(t)\in\mathbb C^{1\times N}$ 为 CU 到 UAV 的信道，$\eta_{i,j}(t)$ 和 $p_j(t)$ 分别为频谱共享指标和 CU 发射功率；$d_i(t)>0$ 是 UAV 到目标的距离，$\xi_0\ge0$、$\rho\ge0$ 为信道参数，$\sigma^2$ 是噪声功率。下面推导最大感知 SINR 对应的接收波束。

任取 UAV $u_i$ 和时隙 $t$。首先将式~\eqref{detail:eq:channel} 代入回波幅值，并利用标量乘法的结合律，逐步得到
\begin{samepage}
\begin{align}
&\left|\bm g_i^H(t)\mathbf A(\theta_i(t))\bm w_i(t)\right|^2\notag\\
&=\left|\sqrt{\xi_0\rho d_i^{-4}(t)}\,
\bm g_i^H(t)\bm a_r(\theta_i(t))\bm a_t^H(\theta_i(t))\bm w_i(t)\right|^2\notag\\
&=\xi_0\rho d_i^{-4}(t)
\left|\bm g_i^H(t)\bm a_r(\theta_i(t))\right|^2
\left|\bm a_t^H(\theta_i(t))\bm w_i(t)\right|^2.\label{detail:eq:factor}
\end{align}
\end{samepage}
将上式代入原始 SINR 定义~\eqref{detail:eq:original-sinr}，得到
\begin{equation}\label{detail:eq:factored}
\frac{\xi_0\rho d_i^{-4}(t)
\left|\bm a_t^H(\theta_i(t))\bm w_i(t)\right|^2
\left|\bm g_i^H(t)\bm a_r(\theta_i(t))\right|^2}
{\displaystyle\sum_{j=1}^J\eta_{i,j}(t)p_j(t)
\left|\bm g_i^H(t)\bm h_{c_j,u_i}^{H}(t)\right|^2+\sigma^2}.
\end{equation}
在上述其他优化变量固定的条件下，式\,(\ref{detail:eq:factored}) 的分子中，下面这一项表示发射波束指向目标的增益及感知信道的路径增益：
\begin{equation}\label{detail:eq:tx-factor}
\xi_0\rho d_i^{-4}(t)
\left|\bm a_t^H(\theta_i(t))\bm w_i(t)\right|^2.
\end{equation}
这一项只包含目标参数、位置及发射波束，没有接收波束 $\bm g_i(t)$；所以改变 $\bm g_i(t)$ 时，式\,(\ref{detail:eq:tx-factor}) 的值不变。若该值大于零，将其作为正因子从最大化问题中提出，不改变最大值对应的接收波束。因此，求解接收波束时只需最大化：
\begin{equation}\label{detail:eq:receive-optimization}
\max_{\|\bm g_i(t)\|^2=1}
\frac{\left|\bm g_i^H(t)\bm a_r(\theta_i(t))\right|^2}
{\displaystyle\sum_{j=1}^J\eta_{i,j}(t)p_j(t)
\left|\bm g_i^H(t)\bm h_{c_j,u_i}^{H}(t)\right|^2+\sigma^2}.
\end{equation}
式\,(\ref{detail:eq:receive-optimization}) 中不再出现 $\bm w_i(t)$，这说明最优接收方向由目标接收导向、CU 信道、匹配、功率及噪声决定，与发射波束的具体取值无关。若式\,(\ref{detail:eq:tx-factor}) 为零，所有接收波束的感知 SINR 均为零。

为求解式\,(\ref{detail:eq:receive-optimization})，先将分母中的干扰和噪声统一写成关于 $\bm g_i(t)$ 的二次型。由于接收波束满足 $\|\bm g_i(t)\|^2=1$，所以
\begin{equation}
\sigma^2=\sigma^2\bm g_i^H(t)\bm g_i(t)
=\bm g_i^H(t)(\sigma^2\mathbf I)\bm g_i(t).
\end{equation}
将噪声和每个 CU 的干扰一起写入括号，原分母等于
\begin{align}\label{detail:eq:denominator}
&\sum_{j=1}^J\eta_{i,j}(t)p_j(t)
\left|\bm g_i^H(t)\bm h_{c_j,u_i}^{H}(t)\right|^2+\sigma^2\notag\\
&=\bm g_i^H(t)\!\left(\sum_{j=1}^J\eta_{i,j}(t)p_j(t)
\bm h_{c_j,u_i}^{H}(t)\bm h_{c_j,u_i}(t)+\sigma^2\mathbf I\right)\!\bm g_i(t).
\end{align}
由于 $\sigma^2>0$ 且 $\eta_{i,j}(t)p_j(t)\ge0$，对任意非零向量 $\bm v\in\mathbb C^{N\times1}$，有
\begin{samepage}
\begin{align}
&\bm v^H\!\left(\sum_{j=1}^J\eta_{i,j}(t)p_j(t)
\bm h_{c_j,u_i}^{H}(t)\bm h_{c_j,u_i}(t)+\sigma^2\mathbf I\right)\!\bm v\notag\\
&=\sum_{j=1}^J\eta_{i,j}(t)p_j(t)
\left|\bm h_{c_j,u_i}(t)\bm v\right|^2+\sigma^2\|\bm v\|^2>0.
\end{align}
\end{samepage}
\begin{samepage}
因此，式\,(\ref{detail:eq:denominator}) 括号中的矩阵为 Hermitian 正定矩阵，其正定平方根及逆平方根均存在。为了把分子和这个分母联系起来，在分子的内积中插入该矩阵的平方根及逆平方根，得到
\begin{align}\label{detail:eq:sqrt-insertion}
&\bm g_i^H(t)\bm a_r(\theta_i(t))\notag\\
&=\left[\left(\sum_{j=1}^J\eta_{i,j}(t)p_j(t)
\bm h_{c_j,u_i}^{H}(t)\bm h_{c_j,u_i}(t)+\sigma^2\mathbf I\right)^{\frac12}\!\bm g_i(t)\right]^H\notag\\
&\quad\times\left[\left(\sum_{j=1}^J\eta_{i,j}(t)p_j(t)
\bm h_{c_j,u_i}^{H}(t)\bm h_{c_j,u_i}(t)+\sigma^2\mathbf I\right)^{-\frac12}\!\bm a_r(\theta_i(t))\right].
\end{align}
\end{samepage}
\begin{samepage}
这里两个矩阵相乘为单位矩阵。对式\,(\ref{detail:eq:sqrt-insertion}) 右侧的两个向量应用复数 Cauchy--Schwarz 不等式 $|\bm x^H\bm y|^2\le\|\bm x\|^2\|\bm y\|^2$，得到
\begin{align}\label{detail:eq:cs-norms}
&\left|\bm g_i^H(t)\bm a_r(\theta_i(t))\right|^2\notag\\
&\le\left\|\left(\sum_{j=1}^J\eta_{i,j}(t)p_j(t)
\bm h_{c_j,u_i}^{H}(t)\bm h_{c_j,u_i}(t)+\sigma^2\mathbf I\right)^{\frac12}\!\bm g_i(t)\right\|^2\notag\\
&\quad\times\left\|\left(\sum_{j=1}^J\eta_{i,j}(t)p_j(t)
\bm h_{c_j,u_i}^{H}(t)\bm h_{c_j,u_i}(t)+\sigma^2\mathbf I\right)^{-\frac12}\!\bm a_r(\theta_i(t))\right\|^2.
\end{align}
\end{samepage}
将两个范数平方分别展开，就得到：
\begin{align}
&\left|\bm g_i^H(t)\bm a_r(\theta_i(t))\right|^2\notag\\
&\le\bm g_i^H(t)\!\left(\sum_{j=1}^J\eta_{i,j}(t)p_j(t)
\bm h_{c_j,u_i}^{H}(t)\bm h_{c_j,u_i}(t)+\sigma^2\mathbf I\right)\!\bm g_i(t)\notag\\
&\quad\times\bm a_r^H(\theta_i(t))\!\left(\sum_{j=1}^J\eta_{i,j}(t)p_j(t)
\bm h_{c_j,u_i}^{H}(t)\bm h_{c_j,u_i}(t)+\sigma^2\mathbf I\right)^{-1}\!\bm a_r(\theta_i(t)).\label{detail:eq:cs}
\end{align}
式\,(\ref{detail:eq:cs}) 右侧第一项恰好是式\,(\ref{detail:eq:denominator}) 的分母；将式\,(\ref{detail:eq:cs}) 两侧除以该正数，就得到式\,(\ref{detail:eq:receive-optimization}) 的上界。Cauchy--Schwarz 不等式取等当且仅当式\,(\ref{detail:eq:sqrt-insertion}) 中的两个向量线性相关；将这一条件两侧乘以对应的逆平方根，可得 $\bm g_i(t)$ 与下式同方向，即相差非零复数倍：
\begin{equation}\label{detail:eq:direction}
\left(\sum_{j=1}^J\eta_{i,j}(t)p_j(t)
\bm h_{c_j,u_i}^{H}(t)\bm h_{c_j,u_i}(t)+\sigma^2\mathbf I\right)^{-1}
\bm a_r(\theta_i(t)).
\end{equation}
由此，满足论文单位范数要求的一个解为
\begin{equation}\label{detail:eq:mvdr}
\bm g_i^{\rm MVDR}(t)=
\frac{\left(\displaystyle\sum_{j=1}^J\eta_{i,j}(t)p_j(t)
\bm h_{c_j,u_i}^{H}(t)\bm h_{c_j,u_i}(t)+\sigma^2\mathbf I\right)^{-1}
\bm a_r(\theta_i(t))}
{\left\|\left(\displaystyle\sum_{j=1}^J\eta_{i,j}(t)p_j(t)
\bm h_{c_j,u_i}^{H}(t)\bm h_{c_j,u_i}(t)+\sigma^2\mathbf I\right)^{-1}
\bm a_r(\theta_i(t))\right\|}.
\end{equation}
经典无失真 MVDR 的问题定义及闭式解证明见附录~\ref{app:classic}；相关文献为 Capon [1] 和 Vorobyov [2，第 2.2 节]。




综上，式~\eqref{detail:eq:mvdr} 即为定理中的式~\eqref{eq:unified-theorem}，也是该时隙最大感知 SINR 的一个单位范数接收波束。



\subsection{第二步：证明 MVDR 接收波束可用于 P0 的全局最优解}

任取 $P0$ 的一个可行解 $\tilde{\bm\zeta}$，保留其中所有时隙的轨迹、频谱匹配、CU 功率、卸载波束、卸载时长及计算频率。对任意 UAV $u_i$ 和时隙 $t$，记干扰加噪声矩阵为
\begin{equation}\label{eq:new-cov}
\mathbf R_i(t)=\sum_{j=1}^{J}\eta_{i,j}(t)p_j(t)\bm h_{c_j,u_i}^{H}(t)\bm h_{c_j,u_i}(t)+\sigma^2\mathbf I.
\end{equation}
由第一步的推导，$\mathbf R_i(t)\succ0$，并且单位范数向量
\begin{equation}\label{eq:new-mvdr}
\bm g_i^{\rm MVDR}(t)=\frac{\mathbf R_i^{-1}(t)\bm a_r(\theta_i(t))}{\|\mathbf R_i^{-1}(t)\bm a_r(\theta_i(t))\|}
\end{equation}
在保持其他变量及感知发射波束不变时，使该 UAV 的感知 SINR 达到最大值。即使原发射目标增益为零，该向量仍为最大 SINR 的一个解。

记原可行解中的感知发射与接收波束分别为 $\tilde{\bm w}_i(t)$ 和 $\tilde{\bm g}_i(t)$，原感知 SINR 为
\begin{equation}\label{eq:new-old-sinr}
\tilde s_i(t)=\frac{\left|\tilde{\bm g}_i^H(t)\mathbf A(\theta_i(t))\tilde{\bm w}_i(t)\right|^2}{\displaystyle\sum_{j=1}^{J}\eta_{i,j}(t)p_j(t)\left|\tilde{\bm g}_i^H(t)\bm h_{c_j,u_i}^{H}(t)\right|^2+\sigma^2}.
\end{equation}
保持 $\tilde{\bm w}_i(t)$ 不变，先以 $\bm g_i^{\rm MVDR}(t)$ 替换原接收波束，记此时的感知 SINR 为 $\hat s_i(t)$。根据上述最大化结论，
\begin{equation}\label{eq:new-ineq}
\hat s_i(t)\ge\tilde s_i(t)\ge0.
\end{equation}
由式~\eqref{eq:new-ineq}，存在 $c_i(t)\in[0,1]$，使得
\begin{equation}\label{eq:new-scale}
c_i^2(t)\hat s_i(t)=\tilde s_i(t).
\end{equation}
这是因为 $c_i^2(t)\hat s_i(t)$ 在 $c_i(t)\in[0,1]$ 上连续取遍 $[0,\hat s_i(t)]$；若 $\hat s_i(t)=0$，则 $\tilde s_i(t)=0$，可取 $c_i(t)=0$。
构造新的感知发射波束 $\bm w'_i(t)=c_i(t)\tilde{\bm w}_i(t)$。因此新解的感知 SINR $s'_i(t)$ 与感知发射功率满足
\begin{samepage}
\begin{align}\label{eq:new-same-sinr}
s'_i(t)&=c_i^2(t)\hat s_i(t)=\tilde s_i(t),\\
\|\bm w'_i(t)\|^2&=c_i^2(t)\|\tilde{\bm w}_i(t)\|^2\le\|\tilde{\bm w}_i(t)\|^2.\label{eq:new-power}
\end{align}
\end{samepage}
对所有 UAV 和时隙分别执行上述构造，所得解记为 $\bm\zeta'$。下面验证其满足 $P0$ 的约束。

首先，接收波束满足单位范数约束；由式~\eqref{eq:new-power} 及原解可行性，感知发射功率仍不超过上限。由于式~\eqref{eq:new-same-sinr} 保持感知 SINR 不变，原有感知 SINR 阈值仍成立。论文模型的感知信息速率为
\begin{equation}\label{eq:new-rate-formula}
R_i^{\rm sen}(t)=\frac{\delta}{2\nu}\log_2\!\left(1+2\sigma_{\rm pre}^2\gamma^2B^3\nu\,s_i(t)\right),
\end{equation}
其中 $s_i(t)$ 为感知 SINR。由此新旧感知速率及数据量相等：
\begin{equation}\label{eq:new-same-rate}
R_i^{\rm sen\prime}(t)=\tilde R_i^{\rm sen}(t),\qquad
\bar D^{\rm sen}R_i^{\rm sen\prime}(t)=\bar D^{\rm sen}\tilde R_i^{\rm sen}(t).
\end{equation}
论文中的感知数据完整卸载和计算时延约束分别为
\begin{align}\label{eq:new-model-off}
D_{u_i}^{\rm off}(t) R_i^{\rm off}(t)&\ge\bar D^{\rm sen}R_i^{\rm sen}(t),\\
\bar D^{\rm sen}+D_{u_i}^{\rm off}(t)+\frac{C^{\rm sen}\bar D^{\rm sen}R_i^{\rm sen}(t)}{f_{u_i}(t)}&\le D_{\max}^{\rm sen}.\label{eq:new-model-delay}
\end{align}
由于卸载波束、卸载时长和计算频率均保持原值，感知任务卸载速率 $R_i^{\rm off}(t)$ 不变；结合原解的可行性，得
\begin{align}
\bar D^{\rm sen}R_i^{\rm sen\prime}(t)&=\bar D^{\rm sen}\tilde R_i^{\rm sen}(t)\le\tilde D_{u_i}^{\rm off}(t)\tilde R_i^{\rm off}(t),\label{eq:new-off}\\
\bar D^{\rm sen}+\tilde D_{u_i}^{\rm off}(t)+\frac{C^{\rm sen}\bar D^{\rm sen}R_i^{\rm sen\prime}(t)}{\tilde f_{u_i}(t)}&\le D_{\max}^{\rm sen}.\label{eq:new-delay}
\end{align}
故感知数据完整卸载与计算时延约束仍满足。原模型中对应的 BS 感知计算能耗也不变，其表达式为
\[
\kappa C^{\rm sen}\bar D^{\rm sen}R_i^{\rm sen}(t) f_{u_i}^2(t).
\]

其次，对于与 UAV $u_i$ 共享频谱的 CU $c_j$，感知阶段的 CU 卸载速率为
\begin{equation}\label{eq:new-cu-rate}
R_{i,j}^{(1)}(t)=B\log_2\!\left(1+\frac{p_j(t)|h_{c_j,\rm BS}(t)|^2}{|\bm h_{u_i,\rm BS}^H(t)\bm w_i(t)|^2+\sigma^2}\right).
\end{equation}
而新解的感知发射干扰满足
\begin{equation}\label{eq:new-interference}
\left|\bm h_{u_i,\mathrm{BS}}^H(t)\bm w'_i(t)\right|^2=c_i^2(t)\left|\bm h_{u_i,\mathrm{BS}}^H(t)\tilde{\bm w}_i(t)\right|^2\le\left|\bm h_{u_i,\mathrm{BS}}^H(t)\tilde{\bm w}_i(t)\right|^2.
\end{equation}
由式~\eqref{eq:new-cu-rate} 和~\eqref{eq:new-interference}，第一阶段速率不降低，即 $R_{i,j}^{(1)\prime}(t)\ge\tilde R_{i,j}^{(1)}(t)$。第二、第三阶段速率及各阶段时长保持不变，而第一阶段在 CU 已卸载数据量中具有非负系数 $\eta_{i,j}(t)\bar D^{\rm sen}$，故 CU 成功卸载约束仍成立。其余变量和约束均未改变，包括跨时隙轨迹约束，故 $\bm\zeta'$ 是 $P0$ 的可行解。

最后比较目标值。令 $E_{P0}(\cdot)$ 表示系统长期平均加权能耗。BS 计算能耗、CU 能耗、UAV 卸载能耗与飞行能耗均不变，感知发射能耗为 $\bar D^{\rm sen}\|\bm w_i(t)\|^2$，故
\begin{align}\label{eq:new-objective}
E_{P0}(\bm\zeta')-E_{P0}(\tilde{\bm\zeta})
&=-\frac{\omega_2\bar D^{\rm sen}}{t_{\max}}\sum_{t=1}^{t_{\max}}\sum_{u_i\in\mathcal U}\bigl(1-c_i^2(t)\bigr)\left\|\tilde{\bm w}_i(t)\right\|^2\le0.
\end{align}
现令 $\bm\zeta^\star$ 为 $P0$ 的一个全局最优解，$E_{P0}^\star=E_{P0}(\bm\zeta^\star)$。对其执行上述构造，得到采用 MVDR 接收波束的可行解 $\bm\zeta^{\rm MVDR}$。由式~\eqref{eq:new-objective}，
\begin{equation}\label{eq:new-opt-upper}
E_{P0}(\bm\zeta^{\rm MVDR})\le E_{P0}(\bm\zeta^\star)=E_{P0}^\star.
\end{equation}
另一方面，由于 $\bm\zeta^{\rm MVDR}$ 对 $P0$ 可行，且 $E_{P0}^\star$ 是全局最小值，有
\begin{equation}\label{eq:new-opt-lower}
E_{P0}(\bm\zeta^{\rm MVDR})\ge E_{P0}^\star.
\end{equation}
结合上述两式可得 $E_{P0}(\bm\zeta^{\rm MVDR})=E_{P0}^\star$。因此，$\bm\zeta^{\rm MVDR}$ 是 $P0$ 的一个全局最优解，定理得证。\hfill$\square$


\section*{参考与可核查来源}
\begin{enumerate}[label={[\arabic*]},leftmargin=*]
\item J. Capon, “High-Resolution Frequency-Wavenumber Spectrum Analysis,” \textit{Proceedings of the IEEE}, vol. 57, no. 8, pp. 1408--1418, 1969. DOI: \href{https://doi.org/10.1109/PROC.1969.7278}{10.1109/PROC.1969.7278}. \href{https://epsc.wustl.edu/~ggeuler/reading/cam_noise_biblio/capon_1969-ieee-high-resolution_frequency-wavenumber_spectrum_analysis.pdf}{原文 PDF}。
\item S. A. Vorobyov, “Principles of Minimum Variance Robust Adaptive Beamforming Design,” \textit{Signal Processing}, vol. 93, no. 12, pp. 3264--3277, 2013. DOI: \href{https://doi.org/10.1016/j.sigpro.2012.10.021}{10.1016/j.sigpro.2012.10.021}. \href{https://users.aalto.fi/~vorobys1/SP13.pdf}{作者 PDF}。
\item 系统模型：多 UAV 辅助 ISCC 系统的感知信道、感知速率、能耗目标和 $P0$ 约束。本文第三节给出基于该模型的构造证明。
\end{enumerate}
\clearpage
\appendix
\section{经典无失真 MVDR 的求解证明}\label{app:classic}
在经典 MVDR 问题中，固定目标接收导向矢量 $\bm a\in\mathbb C^{N\times1}$ 和干扰加噪声协方差矩阵 $\mathbf R$，假设 $\bm a\ne\bm0$、$\mathbf R=\mathbf R^H\succ0$。令接收向量为 $\bm v$，则经典无失真 MVDR 问题为
\begin{align}\label{eq:app-mvdr-problem}
\min_{\bm v\in\mathbb C^{N\times1}}\quad
&\bm v^H\mathbf R\bm v,\\
\text{s.t.}\quad &\bm v^H\bm a=1.\notag
\end{align}
其中目标函数是输出干扰与噪声功率，约束要求目标方向的复数响应为 $1$。这一问题及其闭式解可参见 Capon [1] 和 Vorobyov [2，第 2.2 节]。

\noindent\textbf{证明。}对于任意满足约束 $\bm v^H\bm a=1$ 的向量 $\bm v$，由于 $\mathbf R\succ0$，其正定平方根和逆平方根均存在。将单位矩阵 $\mathbf R^{1/2}\mathbf R^{-1/2}$ 插入内积，并使用 Cauchy--Schwarz 不等式，得到
\begin{align}\label{eq:app-classic-cs}
1&=\left|\bm v^H\bm a\right|^2
=\left|\bigl(\mathbf R^{1/2}\bm v\bigr)^H
\bigl(\mathbf R^{-1/2}\bm a\bigr)\right|^2\notag\\
&\le\bigl(\bm v^H\mathbf R\bm v\bigr)
\bigl(\bm a^H\mathbf R^{-1}\bm a\bigr).
\end{align}
由于 $\bm a\ne\bm0$ 且 $\mathbf R^{-1}\succ0$，有 $\bm a^H\mathbf R^{-1}\bm a>0$，所以任何可行向量都满足
\begin{equation}\label{eq:app-classic-bound}
\bm v^H\mathbf R\bm v\ge
\frac{1}{\bm a^H\mathbf R^{-1}\bm a}.
\end{equation}
令
\begin{equation}\label{eq:app-classic-solution}
\bm v^{\rm MVDR}=
\frac{\mathbf R^{-1}\bm a}{\bm a^H\mathbf R^{-1}\bm a}.
\end{equation}
由于 $\bm a^H\mathbf R^{-1}\bm a$ 为正实数，直接代入可得
\begin{align}\label{eq:app-classic-check}
\bigl(\bm v^{\rm MVDR}\bigr)^H\bm a&=1,\\
\bigl(\bm v^{\rm MVDR}\bigr)^H\mathbf R\bm v^{\rm MVDR}
&=\frac{1}{\bm a^H\mathbf R^{-1}\bm a}.\notag
\end{align}
因此，$\bm v^{\rm MVDR}$ 满足无失真约束并达到式~\eqref{eq:app-classic-bound} 的下界，是问题~\eqref{eq:app-mvdr-problem} 的全局最优解。由 Cauchy--Schwarz 取等条件及无失真约束，这一解唯一。

在本文的单位范数接收约束下，将式~\eqref{eq:app-classic-solution} 的方向归一化，得到
\begin{equation}\label{eq:app-unit-norm}
\bm g^{\rm unit}=\frac{\bm v^{\rm MVDR}}{\|\bm v^{\rm MVDR}\|}
=\frac{\mathbf R^{-1}\bm a}{\|\mathbf R^{-1}\bm a\|}.
\end{equation}
对于非零接收向量，$|\bm v^H\bm a|^2/(\bm v^H\mathbf R\bm v)$ 不受非零复数缩放影响；因此两种归一化对应相同的最大 SINR 方向。经典 MVDR 的 $\bm v^{\rm MVDR}$ 满足无失真约束，本文的 $\bm g^{\rm unit}$ 满足单位范数约束，它们的数值一般不相同。若采用原式分母中的常数噪声项，只有单位范数向量才可直接用 $\bm g^H\mathbf R\bm g$ 表示该分母。\hfill$\square$

\clearpage
\section{MVDR 成立条件及本文模型的对应关系}\label{app:conditions}
本节分别说明固定其他变量时的最大 SINR 结论，以及在 $P0$ 联合优化中可以采用 MVDR 接收波束的条件。后一结论还依赖感知发射波束的配合调整。

\noindent\textbf{一、最大感知 SINR 的条件及本文依据。}
\begin{enumerate}[label=\textbf{(\arabic*)},leftmargin=*]
\item \textbf{秩一目标回波与非零接收导向。}式~\eqref{eq:A} 将感知信道写为接收导向向量与发射导向向量的外积，因而固定其他变量时，感知发射波束只通过标量增益影响 SINR；论文给出的接收导向向量首元素为 $1$，故不为零。若发射目标增益为零，所有接收方向的 SINR 均为零，MVDR 仍是最大值的一个解，但不再是唯一的最大方向。
\item \textbf{干扰加噪声矩阵正定。}原问题中的匹配指标非负、CU 发射功率为正；进一步采用 $\sigma^2>0$ 的正噪声功率假设，则式~\eqref{eq:new-cov} 给出的矩阵为正定矩阵，逆矩阵存在。论文将 $\sigma^2$ 定义为噪声功率，但应用结论时仍需明确其严格为正。
\item \textbf{接收向量的定义域与归一化。}论文本身把波束定义为复向量，并约束 $\|\bm g_i(t)\|=1$；因此可使用式~\eqref{eq:mvdr} 的单位范数形式。原 $P0$ 变量行中三个波束的 $\ge0$ 与复向量定义不一致；根据作者已确认的笔误，应将其理解为复向量域，而不施加逐元素非负实数约束。
\item \textbf{信道与目标方向的可用性。}计算该闭式接收波束需要所用时隙的目标接收导向向量及干扰加噪声矩阵。论文明确假设目标相对方位角已知，信道在时隙内准静态；但准静态并不自动意味着干扰信道被准确估计。若要声称实际计算所得波束达到真实模型的最大 SINR，应另外说明相关信道或协方差和噪声功率已准确获得；有估计误差时，该精确最优性不能直接保证。
\end{enumerate}

\noindent\textbf{二、将 MVDR 用于 $P0$ 最优方案的附加条件及本文依据。}
\begin{enumerate}[label=\textbf{(\arabic*)},leftmargin=*]
\item \textbf{感知发射波束可以缩放。}论文列出的感知发射约束是功率上限，且无已列出的发射功率下限；第三节构造将波束乘以 $[0,1]$ 内的系数，故不会违反该上限。该论证以 $P0$ 已列出完整的波束约束为前提；若另有不能保持于缩放下的硬件或发射约束，需要重新论证。
\item \textbf{保持感知数据量并控制通信干扰。}原模型中的感知速率由感知 SINR 决定；第三节同时调整接收和发射波束，使 SINR 不变，故感知数据量、感知卸载与计算负担不变。感知发射波束缩小不会增大同频 CU 的第一阶段干扰，其余阶段速率和变量不变，因而 CU 的成功卸载约束不恶化。
\item \textbf{能耗权重和最优解存在性。}目标值不增的构造需要 $\omega_2\ge0$。原论文只明确给出权重之和为 $1$，该等式本身不能推出各权重非负；应在模型中补充非负权重假设。第三节还以 $P0$ 存在全局最优解为前提，所得结论是：若最优解存在，则至少有一个采用 MVDR 接收方向。现有证明并未单独证明 $P0$ 一定存在全局最优解。
\end{enumerate}

综上，本文的秩一感知信道、接收单位范数和所列功率约束支持上述推导；正噪声功率、所需信道信息可用、非负能耗权重及 $P0$ 最优解存在性等条件，应按上述区别明确作为假设或进一步验证。

\end{document}
